% =====================================================================
%  Chapter 5 -- Conclusions
%
%  This file is a fragment, not a standalone document. Either paste it
%  into dissertation_chapters.tex immediately before \bibliographystyle,
%  or keep it separate and pull it in with:
%      % =====================================================================
%  Chapter 5 -- Conclusions
%
%  This file is a fragment, not a standalone document. Either paste it
%  into dissertation_chapters.tex immediately before \bibliographystyle,
%  or keep it separate and pull it in with:
%      % =====================================================================
%  Chapter 5 -- Conclusions
%
%  This file is a fragment, not a standalone document. Either paste it
%  into dissertation_chapters.tex immediately before \bibliographystyle,
%  or keep it separate and pull it in with:
%      % =====================================================================
%  Chapter 5 -- Conclusions
%
%  This file is a fragment, not a standalone document. Either paste it
%  into dissertation_chapters.tex immediately before \bibliographystyle,
%  or keep it separate and pull it in with:
%      \input{chapter5_conclusions}
%  It uses the \runin{} heading command defined in that file's preamble.
% =====================================================================

\chapter{Conclusions}

% ---------------------------------------------------------------------
\section{Answering the research question}
% ---------------------------------------------------------------------

This dissertation asked whether vision-language models can judge human skill
level from video, zero-shot, and --- if they fail --- whether that failure is
perceptual or linguistic. The answer to the first part is no. Across four
models, two video-input paradigms, four prompt formats, two viewpoints,
three frame budgets, two trimming conditions and seven activity domains, no
zero-shot configuration produces a reliable, non-degenerate above-chance
result on four-class skill assessment. The one genuine above-chance result
in the project --- Gemini's 76\% on binary climbing, egocentric --- is a
single condition out of roughly 130, surrounded by chance-level neighbours
from the same model on the same task.

The answer to the second part is that the failure is linguistic, and located
at the language-generation step rather than anywhere upstream of it. A
linear probe over frozen CLIP features classifies the same clips at 67.0\%
where every VLM tested scores 24--25\%, so the discriminative signal is
present and linearly accessible in the visual representation before it
reaches a language model. A text-only control given genuine expert
commentary, with no video at all, collapses into the same chance-level band.
Vision alone succeeds; vision-to-language fails; language-to-language fails.
The bottleneck is not what the models see, nor even what they are told, but
what the generation step does with either.

% ---------------------------------------------------------------------
\section{Achievement of the original objectives}
% ---------------------------------------------------------------------

Every one of the eight hypotheses in Table~\ref{tab:hypotheses} was tested as
specified. Three were supported: models do not exceed chance zero-shot (H1);
the failure takes the form of label collapse rather than random guessing,
with 43\% of conditions assigning a single label to 87--100\% of clips (H2);
and a skill-discriminative signal survives in the visual features where
language output fails, recovered by the CLIP probe at 67.0\% (H7).

Five were rejected. An eightfold increase in frame budget does not improve
accuracy (H3). Trimming to the annotated task window helps about as often as
it hurts (H4). The exocentric viewpoint does not consistently beat the
egocentric one, and the direction of the effect reverses between models and
between activities (H5). Performance does not differ meaningfully across
domains, the collapse holding across climbing, dance, basketball, surgical
suturing, music, cooking and soccer (H6). And skill level is not recoverable
from expert language alone (H8) --- the outcome the localisation argument
required, but not the one the hypothesis predicted.

Both outcomes constitute completion of the objectives. A dissertation built
to test hypotheses is not required to have them resolve in one direction,
and these rejections are findings rather than shortfalls. H3 to H6
successively eliminate every input-side explanation for the failure ---
frame budget, irrelevant footage, camera placement, sample size, activity
domain --- and it is that elimination which makes the H7 and H8 controls
interpretable. Had any of the five been supported, the language-bottleneck
claim of §4.7 would not have been available.

% ---------------------------------------------------------------------
\section{Contributions}
% ---------------------------------------------------------------------

\runin{Localisation of the failure to language generation.} The
triangulation of H7 and H8 --- a frozen-CLIP probe succeeding where the VLMs
fail, paired with a text-only control that fails even when handed the
evidence in expert prose --- rules out the visual encoder, rules out the
availability of the evidence in linguistic form, and leaves the conversion
of either into a categorical judgement. This supplies a mechanism rather
than a score, on a task where the discriminative signal is purely dynamic
and evaluative.

\runin{The Video-LLaVA context-ceiling arithmetic.} Sixteen frames at 256
visual tokens each is exactly 4{,}096 tokens --- the entire context window
of a Vicuna-7B backbone --- before a single word of prompt text is added.
This is a general constraint on any fixed-token-per-frame video-language
model with a comparable budget, not a property of these particular results,
and it means ``increase the frame count'' is not a free parameter for such
architectures. The practical danger is that the failure is silent: one
script aborted on an explicit guard, while another proceeded with an
over-length prompt and emitted parseable-looking garbage.

\runin{Paired accuracy-and-distribution reporting.} The JIGSAWS result ---
65.5\% accuracy from a model that answered ``Novice'' on all 29 clips
against a 19/10 imbalanced ground truth --- shows concretely why a bare
accuracy figure is not evidence of discriminative judgement on a
closed-label task. Pairing every accuracy with per-class recall and the
predicted-label distribution is a transferable convention for anyone
evaluating a generative model against a small fixed label set, independent
of skill assessment and of video.

% ---------------------------------------------------------------------
\section{Limitations}
% ---------------------------------------------------------------------

Beyond the caveats recorded in §4.9, three limitations bound these
conclusions at the level of the project as a whole.

The evaluation is zero-shot throughout: no model was fine-tuned and no
labelled exemplars appeared in any prompt, so nothing here speaks to whether
supervised alignment would close the gap the CLIP probe identifies. The
claim is that the zero-shot generation step fails, not that no generation
step could succeed.

Architectural diversity is narrower than four models suggests. Qwen2.5-VL-7B
and Qwen3-VL-8B share a lineage, so the finding is triangulated across three
distinct families --- Qwen, the Vicuna-backed Video-LLaVA, and Gemini ---
rather than four independent systems. That the collapse target
\emph{relocates} across generations rather than diminishing is suggestive,
but two points on one lineage is a thin basis for a claim about generational
trends.

The ground truth is participant-level. EgoExo4D assigns a proficiency label
to a person on the basis of training history and demonstrated competence,
and every take that person recorded inherits it, so a take in which someone
performs atypically is invisible to the scheme. The benchmark therefore
carries some irreducible uncertainty about the clip being judged, not only
about the model judging it --- of which the principled refusals in §4.8,
where Gemini observed that an egocentric clip did not depict the annotated
activity at all, are the visible edge.

% ---------------------------------------------------------------------
\section{Future work}
% ---------------------------------------------------------------------

Two directions follow directly from findings already established here.

The first is supervised alignment. Since a linear probe over frozen features
already recovers 67.0\% from representations the VLMs fail to use, the
natural next experiment trains only the projection layer or the language
head against proficiency labels, leaving the vision encoder untouched. The
prediction implied by §4.7 is specific and falsifiable: because the deficit
is in the mapping rather than in the representation, alignment at the
projection should recover a large fraction of the probe's accuracy at a
small fraction of the cost of end-to-end fine-tuning.

The second is generality. The language-generation bottleneck was established
for one evaluative-judgement task. Whether it holds for others sharing the
same structure --- a visual scene, a continuous underlying quality, and a
categorical verdict the model must commit to, as in aesthetic quality
judgement or visual safety-compliance checking --- or whether it is specific
to the fine-grained kinematic evidence separating a novice climber from an
expert one, is open. The method transfers without modification: the same
three controls apply to any task in that family, and would establish whether
this dissertation has identified a property of skill assessment or a
property of evaluative judgement in vision-language models generally.

%  It uses the \runin{} heading command defined in that file's preamble.
% =====================================================================

\chapter{Conclusions}

% ---------------------------------------------------------------------
\section{Answering the research question}
% ---------------------------------------------------------------------

This dissertation asked whether vision-language models can judge human skill
level from video, zero-shot, and --- if they fail --- whether that failure is
perceptual or linguistic. The answer to the first part is no. Across four
models, two video-input paradigms, four prompt formats, two viewpoints,
three frame budgets, two trimming conditions and seven activity domains, no
zero-shot configuration produces a reliable, non-degenerate above-chance
result on four-class skill assessment. The one genuine above-chance result
in the project --- Gemini's 76\% on binary climbing, egocentric --- is a
single condition out of roughly 130, surrounded by chance-level neighbours
from the same model on the same task.

The answer to the second part is that the failure is linguistic, and located
at the language-generation step rather than anywhere upstream of it. A
linear probe over frozen CLIP features classifies the same clips at 67.0\%
where every VLM tested scores 24--25\%, so the discriminative signal is
present and linearly accessible in the visual representation before it
reaches a language model. A text-only control given genuine expert
commentary, with no video at all, collapses into the same chance-level band.
Vision alone succeeds; vision-to-language fails; language-to-language fails.
The bottleneck is not what the models see, nor even what they are told, but
what the generation step does with either.

% ---------------------------------------------------------------------
\section{Achievement of the original objectives}
% ---------------------------------------------------------------------

Every one of the eight hypotheses in Table~\ref{tab:hypotheses} was tested as
specified. Three were supported: models do not exceed chance zero-shot (H1);
the failure takes the form of label collapse rather than random guessing,
with 43\% of conditions assigning a single label to 87--100\% of clips (H2);
and a skill-discriminative signal survives in the visual features where
language output fails, recovered by the CLIP probe at 67.0\% (H7).

Five were rejected. An eightfold increase in frame budget does not improve
accuracy (H3). Trimming to the annotated task window helps about as often as
it hurts (H4). The exocentric viewpoint does not consistently beat the
egocentric one, and the direction of the effect reverses between models and
between activities (H5). Performance does not differ meaningfully across
domains, the collapse holding across climbing, dance, basketball, surgical
suturing, music, cooking and soccer (H6). And skill level is not recoverable
from expert language alone (H8) --- the outcome the localisation argument
required, but not the one the hypothesis predicted.

Both outcomes constitute completion of the objectives. A dissertation built
to test hypotheses is not required to have them resolve in one direction,
and these rejections are findings rather than shortfalls. H3 to H6
successively eliminate every input-side explanation for the failure ---
frame budget, irrelevant footage, camera placement, sample size, activity
domain --- and it is that elimination which makes the H7 and H8 controls
interpretable. Had any of the five been supported, the language-bottleneck
claim of §4.7 would not have been available.

% ---------------------------------------------------------------------
\section{Contributions}
% ---------------------------------------------------------------------

\runin{Localisation of the failure to language generation.} The
triangulation of H7 and H8 --- a frozen-CLIP probe succeeding where the VLMs
fail, paired with a text-only control that fails even when handed the
evidence in expert prose --- rules out the visual encoder, rules out the
availability of the evidence in linguistic form, and leaves the conversion
of either into a categorical judgement. This supplies a mechanism rather
than a score, on a task where the discriminative signal is purely dynamic
and evaluative.

\runin{The Video-LLaVA context-ceiling arithmetic.} Sixteen frames at 256
visual tokens each is exactly 4{,}096 tokens --- the entire context window
of a Vicuna-7B backbone --- before a single word of prompt text is added.
This is a general constraint on any fixed-token-per-frame video-language
model with a comparable budget, not a property of these particular results,
and it means ``increase the frame count'' is not a free parameter for such
architectures. The practical danger is that the failure is silent: one
script aborted on an explicit guard, while another proceeded with an
over-length prompt and emitted parseable-looking garbage.

\runin{Paired accuracy-and-distribution reporting.} The JIGSAWS result ---
65.5\% accuracy from a model that answered ``Novice'' on all 29 clips
against a 19/10 imbalanced ground truth --- shows concretely why a bare
accuracy figure is not evidence of discriminative judgement on a
closed-label task. Pairing every accuracy with per-class recall and the
predicted-label distribution is a transferable convention for anyone
evaluating a generative model against a small fixed label set, independent
of skill assessment and of video.

% ---------------------------------------------------------------------
\section{Limitations}
% ---------------------------------------------------------------------

Beyond the caveats recorded in §4.9, three limitations bound these
conclusions at the level of the project as a whole.

The evaluation is zero-shot throughout: no model was fine-tuned and no
labelled exemplars appeared in any prompt, so nothing here speaks to whether
supervised alignment would close the gap the CLIP probe identifies. The
claim is that the zero-shot generation step fails, not that no generation
step could succeed.

Architectural diversity is narrower than four models suggests. Qwen2.5-VL-7B
and Qwen3-VL-8B share a lineage, so the finding is triangulated across three
distinct families --- Qwen, the Vicuna-backed Video-LLaVA, and Gemini ---
rather than four independent systems. That the collapse target
\emph{relocates} across generations rather than diminishing is suggestive,
but two points on one lineage is a thin basis for a claim about generational
trends.

The ground truth is participant-level. EgoExo4D assigns a proficiency label
to a person on the basis of training history and demonstrated competence,
and every take that person recorded inherits it, so a take in which someone
performs atypically is invisible to the scheme. The benchmark therefore
carries some irreducible uncertainty about the clip being judged, not only
about the model judging it --- of which the principled refusals in §4.8,
where Gemini observed that an egocentric clip did not depict the annotated
activity at all, are the visible edge.

% ---------------------------------------------------------------------
\section{Future work}
% ---------------------------------------------------------------------

Two directions follow directly from findings already established here.

The first is supervised alignment. Since a linear probe over frozen features
already recovers 67.0\% from representations the VLMs fail to use, the
natural next experiment trains only the projection layer or the language
head against proficiency labels, leaving the vision encoder untouched. The
prediction implied by §4.7 is specific and falsifiable: because the deficit
is in the mapping rather than in the representation, alignment at the
projection should recover a large fraction of the probe's accuracy at a
small fraction of the cost of end-to-end fine-tuning.

The second is generality. The language-generation bottleneck was established
for one evaluative-judgement task. Whether it holds for others sharing the
same structure --- a visual scene, a continuous underlying quality, and a
categorical verdict the model must commit to, as in aesthetic quality
judgement or visual safety-compliance checking --- or whether it is specific
to the fine-grained kinematic evidence separating a novice climber from an
expert one, is open. The method transfers without modification: the same
three controls apply to any task in that family, and would establish whether
this dissertation has identified a property of skill assessment or a
property of evaluative judgement in vision-language models generally.

%  It uses the \runin{} heading command defined in that file's preamble.
% =====================================================================

\chapter{Conclusions}

% ---------------------------------------------------------------------
\section{Answering the research question}
% ---------------------------------------------------------------------

This dissertation asked whether vision-language models can judge human skill
level from video, zero-shot, and --- if they fail --- whether that failure is
perceptual or linguistic. The answer to the first part is no. Across four
models, two video-input paradigms, four prompt formats, two viewpoints,
three frame budgets, two trimming conditions and seven activity domains, no
zero-shot configuration produces a reliable, non-degenerate above-chance
result on four-class skill assessment. The one genuine above-chance result
in the project --- Gemini's 76\% on binary climbing, egocentric --- is a
single condition out of roughly 130, surrounded by chance-level neighbours
from the same model on the same task.

The answer to the second part is that the failure is linguistic, and located
at the language-generation step rather than anywhere upstream of it. A
linear probe over frozen CLIP features classifies the same clips at 67.0\%
where every VLM tested scores 24--25\%, so the discriminative signal is
present and linearly accessible in the visual representation before it
reaches a language model. A text-only control given genuine expert
commentary, with no video at all, collapses into the same chance-level band.
Vision alone succeeds; vision-to-language fails; language-to-language fails.
The bottleneck is not what the models see, nor even what they are told, but
what the generation step does with either.

% ---------------------------------------------------------------------
\section{Achievement of the original objectives}
% ---------------------------------------------------------------------

Every one of the eight hypotheses in Table~\ref{tab:hypotheses} was tested as
specified. Three were supported: models do not exceed chance zero-shot (H1);
the failure takes the form of label collapse rather than random guessing,
with 43\% of conditions assigning a single label to 87--100\% of clips (H2);
and a skill-discriminative signal survives in the visual features where
language output fails, recovered by the CLIP probe at 67.0\% (H7).

Five were rejected. An eightfold increase in frame budget does not improve
accuracy (H3). Trimming to the annotated task window helps about as often as
it hurts (H4). The exocentric viewpoint does not consistently beat the
egocentric one, and the direction of the effect reverses between models and
between activities (H5). Performance does not differ meaningfully across
domains, the collapse holding across climbing, dance, basketball, surgical
suturing, music, cooking and soccer (H6). And skill level is not recoverable
from expert language alone (H8) --- the outcome the localisation argument
required, but not the one the hypothesis predicted.

Both outcomes constitute completion of the objectives. A dissertation built
to test hypotheses is not required to have them resolve in one direction,
and these rejections are findings rather than shortfalls. H3 to H6
successively eliminate every input-side explanation for the failure ---
frame budget, irrelevant footage, camera placement, sample size, activity
domain --- and it is that elimination which makes the H7 and H8 controls
interpretable. Had any of the five been supported, the language-bottleneck
claim of §4.7 would not have been available.

% ---------------------------------------------------------------------
\section{Contributions}
% ---------------------------------------------------------------------

\runin{Localisation of the failure to language generation.} The
triangulation of H7 and H8 --- a frozen-CLIP probe succeeding where the VLMs
fail, paired with a text-only control that fails even when handed the
evidence in expert prose --- rules out the visual encoder, rules out the
availability of the evidence in linguistic form, and leaves the conversion
of either into a categorical judgement. This supplies a mechanism rather
than a score, on a task where the discriminative signal is purely dynamic
and evaluative.

\runin{The Video-LLaVA context-ceiling arithmetic.} Sixteen frames at 256
visual tokens each is exactly 4{,}096 tokens --- the entire context window
of a Vicuna-7B backbone --- before a single word of prompt text is added.
This is a general constraint on any fixed-token-per-frame video-language
model with a comparable budget, not a property of these particular results,
and it means ``increase the frame count'' is not a free parameter for such
architectures. The practical danger is that the failure is silent: one
script aborted on an explicit guard, while another proceeded with an
over-length prompt and emitted parseable-looking garbage.

\runin{Paired accuracy-and-distribution reporting.} The JIGSAWS result ---
65.5\% accuracy from a model that answered ``Novice'' on all 29 clips
against a 19/10 imbalanced ground truth --- shows concretely why a bare
accuracy figure is not evidence of discriminative judgement on a
closed-label task. Pairing every accuracy with per-class recall and the
predicted-label distribution is a transferable convention for anyone
evaluating a generative model against a small fixed label set, independent
of skill assessment and of video.

% ---------------------------------------------------------------------
\section{Limitations}
% ---------------------------------------------------------------------

Beyond the caveats recorded in §4.9, three limitations bound these
conclusions at the level of the project as a whole.

The evaluation is zero-shot throughout: no model was fine-tuned and no
labelled exemplars appeared in any prompt, so nothing here speaks to whether
supervised alignment would close the gap the CLIP probe identifies. The
claim is that the zero-shot generation step fails, not that no generation
step could succeed.

Architectural diversity is narrower than four models suggests. Qwen2.5-VL-7B
and Qwen3-VL-8B share a lineage, so the finding is triangulated across three
distinct families --- Qwen, the Vicuna-backed Video-LLaVA, and Gemini ---
rather than four independent systems. That the collapse target
\emph{relocates} across generations rather than diminishing is suggestive,
but two points on one lineage is a thin basis for a claim about generational
trends.

The ground truth is participant-level. EgoExo4D assigns a proficiency label
to a person on the basis of training history and demonstrated competence,
and every take that person recorded inherits it, so a take in which someone
performs atypically is invisible to the scheme. The benchmark therefore
carries some irreducible uncertainty about the clip being judged, not only
about the model judging it --- of which the principled refusals in §4.8,
where Gemini observed that an egocentric clip did not depict the annotated
activity at all, are the visible edge.

% ---------------------------------------------------------------------
\section{Future work}
% ---------------------------------------------------------------------

Two directions follow directly from findings already established here.

The first is supervised alignment. Since a linear probe over frozen features
already recovers 67.0\% from representations the VLMs fail to use, the
natural next experiment trains only the projection layer or the language
head against proficiency labels, leaving the vision encoder untouched. The
prediction implied by §4.7 is specific and falsifiable: because the deficit
is in the mapping rather than in the representation, alignment at the
projection should recover a large fraction of the probe's accuracy at a
small fraction of the cost of end-to-end fine-tuning.

The second is generality. The language-generation bottleneck was established
for one evaluative-judgement task. Whether it holds for others sharing the
same structure --- a visual scene, a continuous underlying quality, and a
categorical verdict the model must commit to, as in aesthetic quality
judgement or visual safety-compliance checking --- or whether it is specific
to the fine-grained kinematic evidence separating a novice climber from an
expert one, is open. The method transfers without modification: the same
three controls apply to any task in that family, and would establish whether
this dissertation has identified a property of skill assessment or a
property of evaluative judgement in vision-language models generally.

%  It uses the \runin{} heading command defined in that file's preamble.
% =====================================================================

\chapter{Conclusions}

% ---------------------------------------------------------------------
\section{Answering the research question}
% ---------------------------------------------------------------------

This dissertation asked whether vision-language models can judge human skill
level from video, zero-shot, and --- if they fail --- whether that failure is
perceptual or linguistic. The answer to the first part is no. Across four
models, two video-input paradigms, four prompt formats, two viewpoints,
three frame budgets, two trimming conditions and seven activity domains, no
zero-shot configuration produces a reliable, non-degenerate above-chance
result on four-class skill assessment. The one genuine above-chance result
in the project --- Gemini's 76\% on binary climbing, egocentric --- is a
single condition out of roughly 130, surrounded by chance-level neighbours
from the same model on the same task.

The answer to the second part is that the failure is linguistic, and located
at the language-generation step rather than anywhere upstream of it. A
linear probe over frozen CLIP features classifies the same clips at 67.0\%
where every VLM tested scores 24--25\%, so the discriminative signal is
present and linearly accessible in the visual representation before it
reaches a language model. A text-only control given genuine expert
commentary, with no video at all, collapses into the same chance-level band.
Vision alone succeeds; vision-to-language fails; language-to-language fails.
The bottleneck is not what the models see, nor even what they are told, but
what the generation step does with either.

% ---------------------------------------------------------------------
\section{Achievement of the original objectives}
% ---------------------------------------------------------------------

Every one of the eight hypotheses in Table~\ref{tab:hypotheses} was tested as
specified. Three were supported: models do not exceed chance zero-shot (H1);
the failure takes the form of label collapse rather than random guessing,
with 43\% of conditions assigning a single label to 87--100\% of clips (H2);
and a skill-discriminative signal survives in the visual features where
language output fails, recovered by the CLIP probe at 67.0\% (H7).

Five were rejected. An eightfold increase in frame budget does not improve
accuracy (H3). Trimming to the annotated task window helps about as often as
it hurts (H4). The exocentric viewpoint does not consistently beat the
egocentric one, and the direction of the effect reverses between models and
between activities (H5). Performance does not differ meaningfully across
domains, the collapse holding across climbing, dance, basketball, surgical
suturing, music, cooking and soccer (H6). And skill level is not recoverable
from expert language alone (H8) --- the outcome the localisation argument
required, but not the one the hypothesis predicted.

Both outcomes constitute completion of the objectives. A dissertation built
to test hypotheses is not required to have them resolve in one direction,
and these rejections are findings rather than shortfalls. H3 to H6
successively eliminate every input-side explanation for the failure ---
frame budget, irrelevant footage, camera placement, sample size, activity
domain --- and it is that elimination which makes the H7 and H8 controls
interpretable. Had any of the five been supported, the language-bottleneck
claim of §4.7 would not have been available.

% ---------------------------------------------------------------------
\section{Contributions}
% ---------------------------------------------------------------------

\runin{Localisation of the failure to language generation.} The
triangulation of H7 and H8 --- a frozen-CLIP probe succeeding where the VLMs
fail, paired with a text-only control that fails even when handed the
evidence in expert prose --- rules out the visual encoder, rules out the
availability of the evidence in linguistic form, and leaves the conversion
of either into a categorical judgement. This supplies a mechanism rather
than a score, on a task where the discriminative signal is purely dynamic
and evaluative.

\runin{The Video-LLaVA context-ceiling arithmetic.} Sixteen frames at 256
visual tokens each is exactly 4{,}096 tokens --- the entire context window
of a Vicuna-7B backbone --- before a single word of prompt text is added.
This is a general constraint on any fixed-token-per-frame video-language
model with a comparable budget, not a property of these particular results,
and it means ``increase the frame count'' is not a free parameter for such
architectures. The practical danger is that the failure is silent: one
script aborted on an explicit guard, while another proceeded with an
over-length prompt and emitted parseable-looking garbage.

\runin{Paired accuracy-and-distribution reporting.} The JIGSAWS result ---
65.5\% accuracy from a model that answered ``Novice'' on all 29 clips
against a 19/10 imbalanced ground truth --- shows concretely why a bare
accuracy figure is not evidence of discriminative judgement on a
closed-label task. Pairing every accuracy with per-class recall and the
predicted-label distribution is a transferable convention for anyone
evaluating a generative model against a small fixed label set, independent
of skill assessment and of video.

% ---------------------------------------------------------------------
\section{Limitations}
% ---------------------------------------------------------------------

Beyond the caveats recorded in §4.9, three limitations bound these
conclusions at the level of the project as a whole.

The evaluation is zero-shot throughout: no model was fine-tuned and no
labelled exemplars appeared in any prompt, so nothing here speaks to whether
supervised alignment would close the gap the CLIP probe identifies. The
claim is that the zero-shot generation step fails, not that no generation
step could succeed.

Architectural diversity is narrower than four models suggests. Qwen2.5-VL-7B
and Qwen3-VL-8B share a lineage, so the finding is triangulated across three
distinct families --- Qwen, the Vicuna-backed Video-LLaVA, and Gemini ---
rather than four independent systems. That the collapse target
\emph{relocates} across generations rather than diminishing is suggestive,
but two points on one lineage is a thin basis for a claim about generational
trends.

The ground truth is participant-level. EgoExo4D assigns a proficiency label
to a person on the basis of training history and demonstrated competence,
and every take that person recorded inherits it, so a take in which someone
performs atypically is invisible to the scheme. The benchmark therefore
carries some irreducible uncertainty about the clip being judged, not only
about the model judging it --- of which the principled refusals in §4.8,
where Gemini observed that an egocentric clip did not depict the annotated
activity at all, are the visible edge.

% ---------------------------------------------------------------------
\section{Future work}
% ---------------------------------------------------------------------

Two directions follow directly from findings already established here.

The first is supervised alignment. Since a linear probe over frozen features
already recovers 67.0\% from representations the VLMs fail to use, the
natural next experiment trains only the projection layer or the language
head against proficiency labels, leaving the vision encoder untouched. The
prediction implied by §4.7 is specific and falsifiable: because the deficit
is in the mapping rather than in the representation, alignment at the
projection should recover a large fraction of the probe's accuracy at a
small fraction of the cost of end-to-end fine-tuning.

The second is generality. The language-generation bottleneck was established
for one evaluative-judgement task. Whether it holds for others sharing the
same structure --- a visual scene, a continuous underlying quality, and a
categorical verdict the model must commit to, as in aesthetic quality
judgement or visual safety-compliance checking --- or whether it is specific
to the fine-grained kinematic evidence separating a novice climber from an
expert one, is open. The method transfers without modification: the same
three controls apply to any task in that family, and would establish whether
this dissertation has identified a property of skill assessment or a
property of evaluative judgement in vision-language models generally.
